This document presents a commented model for academic works at the Federal University of Ceará, produced with the abntexto-ufc class on top of the abntexto infrastructure. Its purpose is to work simultaneously as a compilable example, a visual guide, and a validation corpus by explaining the organization of pre-textual, textual, and post-textual elements and by distinguishing normative requirements, UFC institutional guidance, technical template policies, and editorial examples. The reference base includes the current editions adopted by the project for academic-work presentation, citations, references, abstracts, progressive section numbering, table of contents, research projects, indexes, spines, and numerical tables. The textual body demonstrates page size and margins, typography, spacing, paragraphs, pagination, sections, footnotes, citations, and references. It also presents figures, charts, framed tables, numerical tables, source code, algorithms, equations, itemized subdivisions, glossary, appendices, annexes, and index, indicating whether each relevant characteristic comes from a standard, an institutional rule, or a template decision. Repository validation combines semantic and geometric PDF inspections, profile and engine matrices, compatibility with an Overleaf-like environment, literal-font certification on Windows, and PDF/A validation. The model reduces manual formatting work but does not replace academic review or the verification of current requirements established by the applicable course, graduate program, and institutional deposit procedure.

\keywords{academic works; ABNT; LaTeX; standardization; UFC.}
